\section{Conclusion}\label{sec:conclusion}
Exact sparse network--simplex formulations are controlled by unobserved
circulation freedom. The observation-sensitive construction retains one
coordinate per independent unobserved cycle for each locally observed
block--label pair, so it gives a direct hull description when the unobserved
subgraphs are forests and a smaller exact extension otherwise. The count
describes this construction and the associated completion by individual
products, not unrestricted extension complexity.

Eliminating the remaining state coordinates raises a different question:
which coefficients are needed in the original coordinates? Bounded block
cycle rank permits finite support and recovery libraries and a uniform
rank-dependent coefficient bound. Long flat chains admit a second positive
result: unit flow/product coefficients suffice for every observation pattern
with at most three observed labels, and four labels can force a ratio of
two. The rank-three $K_4$ example, the sparse Fibonacci constructions on
simple series--parallel graphs of maximum degree three, and the two-label
universality transfer show that neither small treewidth alone nor a small
number of labels alone implies a uniform coefficient bound over general
networks. These are obstructions to simple projected descriptions; exact
extended formulations and polynomial-time separation remain available.

For modeling, the results offer a choice between retaining a small set of
state coordinates and generating complete structural cuts. For exact
verification, the recovery constructions turn membership into explicit
state flows and convex decompositions, including on degenerate boundaries.
The experiments show substantial formulation-size savings, particularly when
labels occur in different blocks, but global merging of unused labels and
native flow bounds often make a conventional LP the fastest method. A runtime
benefit must be demonstrated for the intended workload, not inferred from a
smaller formulation or an operation bound.

The results concern the equality-flow component hull with a shared simplex
and selected products. Side constraints may be added to this hull to obtain
a valid relaxation, but exactness for a larger coupled model requires a
separate argument.
